\documentclass[12pt]{article}
\usepackage[utf8]{inputenc}
\usepackage[T1]{fontenc}
\usepackage{hyperref}
\title{Yimas Language Resources for K-12 Curriculum Development}
\author{Compiled by Claude Sonnet 5.5 (claude-sonnet-5-5) for Andr\'as Kornai}
\date{}
\begin{document}
\maketitle

\section{Introduction}
Yimas (ISO 639-3 \texttt{yee}, Glottocode yima1243) is a Lower Sepik language of the Sepik basin, Papua New Guinea. It is polysynthetic and morphologically very complex. The prompt estimates about 300 speakers. One search result instead gave about 50 speakers and said no speakers are under 30. The figures disagree and I could not resolve them. The searches turned up scholarly descriptions only. They found no school materials, primers or dictionaries.

\section{Primary Educational Resources}
The searches found no primers, textbooks or school programmes for Yimas. Any curriculum would have to be built from the linguistic descriptions below, together with community input.

\section{Linguistic and Reference Materials}
\begin{itemize}
\item Wikipedia overview of the Yimas language: \url{https://en.wikipedia.org/wiki/Yimas_language}
\item Grammatical sketch of Yimas in the Cambridge Handbook of Role and Reference Grammar: \url{https://resolve.cambridge.org/core/books/cambridge-handbook-of-role-and-reference-grammar/grammatical-sketch-of-yimas-lower-sepik-papua-new-guinea/F9948865D85C8428AEF8BD32B643106B}
\item Foley (2005), listed on the Trans New Guinea database: \url{https://transnewguinea.org/source/foley2005?page=8}
\item Foley (1986), PDF on Papuan languages: \url{https://cdstar.eva.mpg.de/bitstreams/EAEA0-72DE-267C-1377-0/Foley1986.pdf}
\item University of Oregon thesis on serial verb constructions in Yimas (Mathieu-Reeves, 2009): \url{https://scholarsbank.uoregon.edu/xmlui/handle/1794/9836?show=full}
\item Yimas numeral system data: \url{https://lingweb.eva.mpg.de/channumerals/Yimas.htm}
\item APiCS survey chapter on Yimas-Arafundi Pidgin, a related contact variety: \url{https://apics-online.info/surveys/69}
\end{itemize}

\section{Religious and Community Resources}
None were found in the searches.

\section{Digital and Audio Resources}
No audio, video or software resources were found in the searches.

\section{Curriculum Development Context}
The available material is scholarly. A dictionary was not found. Any K-12 curriculum would need new work with Yimas speakers: documenting vocabulary, producing readers, and writing an agreed orthography. The academic grammars can support that work.
\end{document}